\section{Анализ предметной области}

\subsection{Характеристика деятельности редакционно-издательского отдела вуза}

Редакционно-издательский отдел (РИО) высшего учебного заведения является структурным подразделением, обеспечивающим подготовку, проверку, оформление и выпуск учебно-методической, научной и справочной литературы, используемой в образовательном процессе. К числу основных материалов, поступающих в РИО, относятся методические указания, учебные пособия, практикумы, монографии, сборники научных трудов и иные издания.

Современный образовательный процесс требует постоянного обновления учебно-методической базы. Изменение федеральных государственных образовательных стандартов, развитие научных направлений, внедрение новых технологий обучения и цифровых платформ обуславливают регулярную подготовку новых материалов преподавателями университета.

РИО выполняет функции организационного и нормативного сопровождения издательской деятельности вуза. Сотрудники подразделения принимают материалы от авторов, проверяют их соответствие установленным требованиям оформления, контролируют наличие обязательных структурных элементов, подготавливают замечания, формируют титульные листы, а также ведут учёт поступивших и утверждённых работ.

В условиях большого количества авторов и поступающих документов особую значимость приобретают скорость обработки материалов, точность проверки и удобство взаимодействия между преподавателями и сотрудниками отдела. Эффективность деятельности РИО напрямую влияет на своевременное обеспечение образовательного процесса актуальными учебными материалами.

\subsection{Проблематика традиционного подхода к организации проверки материалов}

Традиционный подход к работе редакционно-издательского отдела, основанный на ручной обработке документов, сопровождается рядом существенных недостатков, снижающих эффективность работы подразделения:
\begin{enumerate}
\item{Высокая трудоёмкость проверки оформления.} В большинстве случаев сотрудники РИО вручную проверяют поступившие документы на соответствие установленным требованиям: параметры шрифта, интервалы, отступы, структура заголовков, оформление списков литературы, таблиц, рисунков и других элементов. При большом количестве работ такой процесс занимает значительное время.

\item{Высокая вероятность человеческих ошибок.} Ручная проверка не гарантирует обнаружение всех нарушений оформления. Из-за большого объёма текста, однотипности операций и человеческого фактора часть ошибок может остаться незамеченной, что приводит к возврату материалов на доработку уже на поздних этапах.

\item{Длительное взаимодействие с авторами.} После выявления замечаний сотрудник вынужден вручную составлять перечень ошибок либо вносить исправления в режиме рецензирования документа. Это увеличивает сроки согласования и повторной подачи материалов.

\item{Отсутствие централизованного хранилища работ.} Готовые и поступившие материалы часто хранятся в отдельных папках, на персональных компьютерах или сетевых дисках. Это затрудняет поиск ранее утверждённых изданий, ведение архива и получение статистических сведений.

\item{Сложность формирования отчётности.} Для подготовки отчётов по количеству проверенных работ, числу авторов, срокам обработки и видам материалов требуется вручную собирать данные из различных источников. Это увеличивает нагрузку на сотрудников и снижает оперативность получения информации.

\item{Отсутствие унифицированных шаблонов.} Авторы нередко подготавливают титульные листы и иные элементы документа самостоятельно, что приводит к многочисленным отклонениям от утверждённых образцов и повторным исправлениям.
\end{enumerate}

Таким образом, ручной подход к организации работы РИО не обеспечивает необходимой скорости, точности и удобства обработки документов в современных условиях.

\subsection{Обоснование необходимости автоматизации}

Решение перечисленных проблем возможно путём внедрения специализированной информационно-аналитической системы для автоматизации деятельности редакционно-издательского отдела вуза.

Разрабатываемая система должна обеспечить:

\begin{itemize}
	\item загрузку пользователями учебно-методических и научных материалов через веб-интерфейс;
	\item автоматическую проверку документов на соответствие требованиям оформления;
	\item формирование подробного отчёта с перечнем найденных ошибок;
	\item автоматическое создание файла формата DOCX с выделенными нарушениями;
	\item автоматическое сохранение корректно оформленных работ в централизованном хранилище;
	\item ведение статистики поступивших, проверенных и утверждённых материалов;
	\item генерацию титульных листов по утверждённым шаблонам;
	\item разграничение прав доступа между авторами, редакторами и администраторами системы;
	\item удобный интерфейс взаимодействия пользователей с системой.
\end{itemize}

Автоматизация позволит существенно сократить время проверки материалов, снизить количество ошибок, ускорить процесс согласования документов и обеспечить прозрачный контроль деятельности подразделения.

\subsection{Анализ существующих систем автоматизации документооборота и проверки документов}

На современном рынке представлены различные программные решения, связанные с управлением документами, совместной работой и проверкой текстовых материалов. Однако большинство из них не ориентировано на специфику работы редакционно-издательского отдела вуза.

\subsubsection{Microsoft Word и встроенные средства рецензирования}

Текстовый редактор Microsoft Word широко применяется для подготовки и проверки документов. Он предоставляет следующие возможности:

\begin{itemize}
	\item создание и редактирование документов;
	\item проверка орфографии и грамматики;
	\item режим рецензирования и внесения исправлений;
	\item использование шаблонов оформления.
\end{itemize}

Недостатком данного подхода является отсутствие автоматизированной проверки документа на соответствие локальным требованиям конкретного вуза. Проверка структуры, параметров оформления и соответствия внутренним стандартам выполняется вручную.

\subsubsection{Google Docs}

Google Docs является облачным сервисом для совместной работы с документами. Его возможности включают:

\begin{itemize}
	\item онлайн-редактирование файлов;
	\item совместную работу нескольких пользователей;
	\item систему комментариев и предложений;
	\item хранение документов в облаке.
\end{itemize}

Несмотря на удобство совместного редактирования, сервис не предоставляет специализированных механизмов автоматической проверки учебно-методических материалов по требованиям университета, а также не содержит функций аналитической отчётности для РИО.

\subsubsection{Системы электронного документооборота}

Корпоративные системы электронного документооборота предназначены для регистрации, маршрутизации и хранения документов. Они позволяют:

\begin{itemize}
	\item вести архив документов;
	\item организовывать согласование файлов;
	\item разграничивать права доступа;
	\item контролировать исполнение задач.
\end{itemize}

Однако подобные системы являются универсальными решениями и не включают инструменты интеллектуального анализа оформления документов DOCX, автоматического подчёркивания ошибок и генерации титульных листов.

\subsubsection{Вывод по анализу аналогов}

Проведённый анализ существующих решений показал, что универсальные офисные продукты и системы электронного документооборота способны решать лишь отдельные задачи редакционно-издательского отдела.

Одни системы обеспечивают редактирование документов, другие --- хранение файлов и организацию согласования, однако ни одно из рассмотренных решений не сочетает в себе функции автоматической проверки оформления методических материалов, генерации отчётов об ошибках, формирования исправленного DOCX-файла и накопления статистики работы подразделения.

Следовательно, наиболее целесообразным является создание собственной специализированной информационно-аналитической системы, учитывающей особенности документооборота и нормативные требования конкретного высшего учебного заведения.

\subsection{Постановка задач на разработку}

На основании проведённого анализа были сформулированы следующие задачи, которые должна решать разрабатываемая система:

\begin{enumerate}
	\item Обеспечить регистрацию и авторизацию пользователей с разграничением прав доступа.
	\item Реализовать загрузку методических, учебных и научных материалов в формате DOCX.
	\item Выполнить автоматическую проверку документов на соответствие установленным требованиям оформления.
	\item Формировать отчёт с перечнем найденных нарушений.
	\item Создавать копию документа DOCX с визуальным выделением ошибок.
	\item Автоматически сохранять корректно оформленные работы в централизованной базе данных.
	\item Реализовать хранение и поиск ранее загруженных материалов.
	\item Предоставить средства формирования статистической и аналитической отчётности.
	\item Реализовать модуль автоматического создания титульных листов по шаблону вуза.
	\item Создать удобный веб-интерфейс для пользователей и административную панель для управления системой.
\end{enumerate}